% One-slide talk. Self-contained LaTeX source.
% Compile oslide.tex with XeLaTeX or Tectonic. Exactly one page, 16:9.
\documentclass[17pt]{extarticle}
\usepackage[paperwidth=32cm,paperheight=18cm,margin=0cm]{geometry}
\usepackage{fontspec}
\setmainfont{cmunss.otf}[BoldFont=cmunsx.otf,ItalicFont=cmunsi.otf,BoldItalicFont=cmunso.otf]
\setsansfont{cmunss.otf}[BoldFont=cmunsx.otf,ItalicFont=cmunsi.otf,BoldItalicFont=cmunso.otf]
\usepackage{polyglossia}
\setdefaultlanguage{russian}
\usepackage{amsmath}
\usepackage{tikz}
\usetikzlibrary{arrows.meta}
\pagestyle{empty}
\setlength{\parindent}{0pt}
\renewcommand{\familydefault}{\sfdefault}
\newcommand{\Pari}{P_{\mathrm{ari\_max}}}
\newcommand{\Psilh}{P_{\mathrm{silh\_max}}}
\definecolor{ink}{HTML}{142D4E}
\definecolor{accent}{HTML}{176C8C}
\definecolor{muted}{HTML}{627080}
\definecolor{grid}{HTML}{DFE5EA}
\begin{document}
\null
\begin{tikzpicture}[remember picture,overlay,x=1cm,y=1cm]
\begin{scope}[shift={(current page.south west)}]
\node[anchor=north west,text=ink,text width=30cm,font=\fontsize{25}{29}\selectfont\bfseries] at (1,17.15)
 {Выбор кластерного разбиения по ограниченной внешней информации};
\node[anchor=north west,font=\fontsize{21}{26}\selectfont,text=muted] at (1,15.65)
 {Идея: подключать геометрию, когда оценки по меткам близки с учётом их разброса};

% Horizontal axis: gain in held-out ARI; every circle is one annotation repeat.
% Plot coordinate x = 15.0 + 48 * delta; a shared scale is used for all rows.
\foreach \v/\label in {-0.10/{-$0{,}10$},-0.05/{-$0{,}05$},0/{$0$},0.05/{+$0{,}05$},0.10/{+$0{,}10$}} {
 \pgfmathsetmacro{\xx}{15+48*\v}
 \draw[grid,line width=0.6pt] (\xx,6.45)--(\xx,12.9);
 \node[anchor=north,font=\fontsize{17}{20}\selectfont,text=muted] at (\xx,6.25) {\label};
}
\draw[muted,dashed,line width=1pt] (15,6.45)--(15,12.9);
\node[anchor=east,align=right,text width=6.7cm,font=\fontsize{18}{22}\selectfont] at (7.95,12.0)
 {Случайно, 20 меток};
\node[anchor=east,align=right,text width=6.7cm,font=\fontsize{18}{22}\selectfont] at (7.95,10.4)
 {Случайно, 50 меток};
\node[anchor=east,align=right,text width=6.7cm,font=\fontsize{18}{22}\selectfont] at (7.95,8.8)
 {5 классов, 20 меток};
\node[anchor=east,align=right,text width=6.7cm,font=\fontsize{18}{22}\selectfont] at (7.95,7.2)
 {5 классов, 50 меток};
\fill[muted!65] (15.000000,11.760000) circle[radius=0.068cm];
\fill[muted!65] (15.000000,11.880000) circle[radius=0.068cm];
\fill[muted!65] (19.582444,12.000000) circle[radius=0.068cm];
\fill[muted!65] (15.000000,12.120000) circle[radius=0.068cm];
\fill[muted!65] (17.666402,12.240000) circle[radius=0.068cm];
\fill[accent] (16.449769,12.130000)--(16.579769,12.000000)--(16.449769,11.870000)--(16.319769,12.000000)--cycle;
\node[anchor=south west,text=accent,font=\fontsize{17}{20}\selectfont\bfseries] at (16.579769,12.150000) {$ +0{,}030 $};
\fill[muted!65] (15.000000,10.160000) circle[radius=0.068cm];
\fill[muted!65] (15.000000,10.280000) circle[radius=0.068cm];
\fill[muted!65] (15.000000,10.400000) circle[radius=0.068cm];
\fill[muted!65] (14.017889,10.520000) circle[radius=0.068cm];
\fill[muted!65] (14.017889,10.640000) circle[radius=0.068cm];
\fill[accent] (14.607156,10.530000)--(14.737156,10.400000)--(14.607156,10.270000)--(14.477156,10.400000)--cycle;
\node[anchor=south west,text=accent,font=\fontsize{17}{20}\selectfont\bfseries] at (14.737156,10.550000) {$ -0{,}008 $};
\fill[muted!65] (15.000000,8.560000) circle[radius=0.068cm];
\fill[muted!65] (15.000000,8.680000) circle[radius=0.068cm];
\fill[muted!65] (17.666402,8.800000) circle[radius=0.068cm];
\fill[muted!65] (15.000000,8.920000) circle[radius=0.068cm];
\fill[muted!65] (9.435445,9.040000) circle[radius=0.068cm];
\fill[accent] (14.420369,8.930000)--(14.550369,8.800000)--(14.420369,8.670000)--(14.290369,8.800000)--cycle;
\node[anchor=south west,text=accent,font=\fontsize{17}{20}\selectfont\bfseries] at (14.550369,8.950000) {$ -0{,}012 $};
\fill[muted!65] (15.000000,6.960000) circle[radius=0.068cm];
\fill[muted!65] (15.000000,7.080000) circle[radius=0.068cm];
\fill[muted!65] (15.000000,7.200000) circle[radius=0.068cm];
\fill[muted!65] (15.000000,7.320000) circle[radius=0.068cm];
\fill[muted!65] (15.000000,7.440000) circle[radius=0.068cm];
\fill[accent] (15.000000,7.330000)--(15.130000,7.200000)--(15.000000,7.070000)--(14.870000,7.200000)--cycle;
\node[anchor=south west,text=accent,font=\fontsize{17}{20}\selectfont\bfseries] at (15.130000,7.350000) {$ +0{,}000 $};
\node[anchor=north,font=\fontsize{19}{23}\selectfont,text=ink] at (15,5.35)
 {$\Delta\mathrm{ARI}=\mathrm{ARI}_{\rm test}(\Psilh)-\mathrm{ARI}_{\rm test}(\Pari)$};
\node[anchor=north,font=\fontsize{16}{19}\selectfont,text=muted] at (15,4.5)
 {левее нуля --- хуже; правее --- лучше};

\node[anchor=north west,text=ink,font=\fontsize{22}{26}\selectfont\bfseries] at (23.6,12.8)
 {Разбиения};
\node[anchor=north west,text width=7.3cm,align=left,font=\fontsize{18}{22}\selectfont,text=ink] at (23.6,11.45)
 {$\Pari$: максимум ARI на доступных метках.};
\node[anchor=north west,text width=7.3cm,align=left,font=\fontsize{18}{22}\selectfont,text=ink] at (23.6,9.2)
 {$\Psilh$: максимум silhouette среди близких по ARI вариантов.\par\vspace{0.3cm}
 Допуск --- по бутстрепному разбросу разности ARI.};
\node[anchor=north west,text width=14cm,align=left,font=\fontsize{14}{17}\selectfont,text=muted] at (8.4,13.9)
 {Круг --- повтор; ромб --- среднее.};

\node[anchor=west,font=\fontsize{23}{28}\selectfont\bfseries,text=accent] at (1,3.12)
 {20 случайных меток: $+0{,}030$ ARI};
\node[anchor=west,font=\fontsize{21}{25}\selectfont,text=ink] at (18.6,3.12)
 {При 50 метках: $-0{,}008$};
\node[anchor=west,font=\fontsize{20}{24}\selectfont,text=ink] at (1,2.15)
 {Вопрос работы: когда геометрия помогает и как это зависит от покрытия классов?};
\node[anchor=south west,text width=30cm,font=\fontsize{15}{18}\selectfont,text=muted] at (1,0.52)
 {Пилот Digits: 600 объектов, 10 классов, 8 готовых разбиений; ARI на 300 скрытых проверочных метках.\\
 5 повторений разметки на условие, 30 бутстрепов.};
\end{scope}
\end{tikzpicture}
\end{document}
